\section{Conclusion}
\label{sec:conclusion}

We introduced DeXposure-FM, a time-series graph foundation model for measuring and forecasting inter-protocol credit exposures in decentralised finance. Trained on the DeXposure dataset spanning over 4,300 protocols across 602 blockchains from 2020 to 2025, the model learns joint representations of protocol-level dynamics, edge-level exposures, and sector structure. We evaluated DeXposure-FM on three tasks: multi-step forecasting, financial stability analysis during historical stress events, and imputation of missing edges. GraphPFN-based variants maintain stable predictive performance (AUROC $>$0.91) across forecast horizons up to 14 weeks, outperforming temporal GNN baselines that degrade at longer horizons. The model prioritises economically significant edges, achieving substantially lower TVL-weighted prediction error than alternatives.

Beyond forecasting accuracy, DeXposure-FM provides a measurement layer for macroprudential monitoring: systemic importance scores, sector spillover matrices, and early-warning indicators derived from network structure. These outputs can support regulators in identifying concentration risks and conducting stress tests on DeFi infrastructure.

The current work is subject to several limitations. The model observes only on-chain DeFi data and cannot capture exposures involving centralised exchanges, off-chain reserves, or traditional financial institutions. Weekly aggregation smooths out high-frequency dynamics. Future research could integrate off-chain data sources, refine exposure measures with asset-level risk weights, and extend the framework to real-time monitoring at finer temporal resolution.
